%قبل از شروع مسئله ترجیحاً چیزی ننویسید. اما اگر خواستید بنویسید، در قالب لیستی (فقرات و شمارش) بنویسید.
%برای اضافه کردن نمره، بعدی از شروع‌مسئله، [x] را اضافه کنید و به جای x  نمره مورد نظر را قرار دهید. برای زیرمسئله‌ها هم به همین صورت است.
%\شروع‌مسئله[10]
\شروع‌مسئله[24] \متن‌سیاه{تحلیل مجانبی}. در هر قسمت علاوه بر پاسخ درست، روابطی که استفاده کردید را نوشته و به طور مختصر توضیح دهید.
در سؤال‌های زیر مبنای $\log$ برابر با ۲ می‌باشد.
\شروع{زیرمسائل}
\زیرمسئله[6]
با استفاده از قضیۀ اصلی، رابطۀ زیر را حل کنید.
\begin{center}
	$T(n) = 49 T(n/7) + n^2 (\log n)^3$
\end{center}

\پاسخ{
	طبق حالت دوم قضیه اصلی داریم:
	
	$a=\log_{7}49 \rightarrow T(n)=O(n^2 log^4) n$
	
}
\نمره‌دهی{
	گرفتن لگاریتم مذکور 2 نمره و نوشتن حالت قضیه اصلی و حل کامل مسئله 4 نمره. 
	
}


\زیرمسئله[6] رابطۀ بازگشتی زیر را حل کنید.
\begin{center}
$T(n)=T(\sqrt{n}) + \log n$
\end{center}
%%2011FQ1

\پاسخ{
	با دادن تغییر متغیر $n = 2^m$ و $T(n) = S(m)$ مسئله به حالت اول قضیه اصلی تبدیل می شود و داریم:
	
	$S(m) = O(m) \rightarrow T(n) = O(\log n)$
	
	لازم به ذکر است مسئله بعد از تغییر متغیر با 
	استفاده از درخت بازگشتی هم قابل حل است.
	
}
\نمره‌دهی{
	
	
	تغییر متغیر دو نمره و ادامه راه و رسیدن به جواب درست 4 نمره. 
	
}



\زیرمسئله[6]
الگوریتمی داریم که زمان اجرای آن
$O(n^{2-(3/2)\phi} + n^{2\beta} + n^{(1-\beta)(2-\phi)+1.5\beta})$
است. پارامترهای $\beta$ و $\phi$ را می‌توانیم به دلخواه از اعداد بین صفر تا یک انتخاب کنیم. هدف این که با انتخاب پارامترهای مناسب زمان اجرای الگوریتم را تا حد ممکن کاهش دهیم. پارامترها مناسب و زمان اجرای الگوریتم با این پارامترها چقدر خواهد بود؟ (راهنمایی: زمانی الگوریتم کمترین زمان را دارد که توان سه جمله یکسان باشد.)

\پاسخ{
$\phi = \frac{5}{4} - \frac{\sqrt{\frac{43}{3}}}{4} \quad \beta = \frac{1}{16} - \frac{\sqrt{129}}{16}$

طبق راهنمایی سوال کافیست دو معادله دو مجهول زیر را حل کنیم.

$$2 - 1.5\phi = 2\beta$$
$$2\beta = (1-\beta)(2-\phi) + 1.5\beta$$

این معادلات دو مجموعه جواب دارند اما فقط یک مجموعه آن‌ها معتبر (در بازه ۰ و ۱) است.

}

\نمره‌دهی{
راه حل با جواب کاملا اشتباه (بدون شباهت به جواب اصلی) که فقط به معادلات درست اشاره کرده باشد ۱ نمره گرفته و سایر راه حل‌ها به تناسب نمره می‌گیرند.

}

\زیرمسئله[6] رابطۀ بازگشتی زیر را حل کنید. 
\begin{center}
	$T(n)=T(n / \log n)+O(1)$
\end{center}

\پاسخ{
	$O(\log n / \log\log n)$
	
	می‌توان پاسخ بالا را به کمک درخت بازگشتی حدس زد. برای اثبات آن، تغییر متغیر
	$$s = \log n \quad f(s) = f(s - \log s) + O(1)$$
	را در نظر بگیرید. می‌خواهیم نشان دهیم
	$$f(s) = O(\frac{s}{\log s})$$
	برای اثبات قضیه فوق ابتدا لم زیر را ثابت می‌کنیم
	$$\exists A,x_0 \forall x>x_0 \quad A\log x \leq \log(x - \log x)$$
	برای اثبات لم فوق دقت کنید که
	$$\log(x - \log x) = \log(x) + \log(1 - \frac{\log x}{x})$$
	عبارت
	$log(1 - \frac{\log x}{x})$
	هنگامی که
	$x \to \infty$
	به صفر میل می‌کند بنابراین به ازای 
	$\epsilon$
	به اندازه دلخواه کوچک نامنفی داریم
	$$\exists x_0 \forall x > x_0\quad \log(1 - \frac{\log x}{x}) \geq -\epsilon \implies \log(x - \log x) \geq \log x - \epsilon \geq \log x(1 - \frac{\epsilon}{\log x_0})$$

	پس کافیست داشته باشیم
	$A < 1 - \frac{\epsilon}{\log x_0}$.
	حال اگر دنباله
	$t_n$
	را به شکل
	$t_{n+1} = t_n - \log t_n$
	تعریف کنیم آنگاه
	$f(n) \in O(n^*)$
	به گونه‌ای که
	$t_{n^*} = 0$.
	$$\log t_{n+1} \geq A \log t_n \implies t_n^* \geq s - \log s\sum_{i=1}^{n}A^i \implies \frac{1 - A^{n^*}}{1 - A} \leq \frac{s}{\log s}$$
	تعریف می‌کنیم
	$A = 1 - \epsilon$
	.آنگاه
	$$1 - A^{n^*}  = 1 - (1 - \epsilon)^{n^*} \leq 1 - (1 - \frac{\epsilon n^*}{2}) \implies \frac{\epsilon n^*}{2} \leq \frac{\epsilon s}{\log s}$$
	$$n^* \in O(\frac{s}{\log s}) \implies f(s) \in O(\frac{s}{\log s}) \implies T(n) \in O(\frac{\log n}{\log \log n})$$
	
}

\نمره‌دهی{
	در این سوال به پاسخ‌هایی که با استفاده از درخت بازگشتی به جواب رسیده‌اند نیز اگر استدلال درست باشد نمره کامل و برای استدلال‌های نادرست نیز به تناسب راه حل نمره در نظر گرفته شده است.
}

%اگر مسئله چند قسمت داشت به این صورت عمل کنید:
%\شروع{زیرمسائل}
%\زیرمسئله[0] % مسئلۀ ۱آ
%\زیرمسئله % مسئلۀ ۱ب
%\زیرمسئله % مسئلۀ ۱پ
%\زیرمسئله % مسئلۀ ۱ت
\پایان{زیرمسائل}




%%%%%%%%%%%%%%%%%%%%%%%%%%%%%%%%%%%%%%%%%%%%%%%%%%%%%
% برخی دستورات مفید برای تمرین‌ها. %
%%%%%%%%%%%%%%%%%%%%%%%%%%%%%%%%%%%%%%%%%%%%%%%%%%%%%

% فرمول‌های نمونه:
%$f(x) = \Theta(x)$
%$T(x, y) \leq \log(x) + 2^y + \binom{2n}{n}$
% {\tt code\_function}


% لیست موارد بدون شماره‌گذاری:
%\شروع{فقرات}
%    \فقره مورد اول 
%        \شروع{فقرات}
%            \فقره مورد اول داخلی 
%                \شروع{فقرات}
%                    \فقره اولین مورد خیلی داخلی 
%                    \فقره دومین مورد خیلی داخلی 
%                \پایان{فقرات}
%            \فقره مورد دوم داخلی 
%        \پایان{فقرات}
%    \فقره مورد دوم 
%\پایان{فقرات}

% لیست موارد شماره‌گذاری شده:
%\شروع{شمارش}
%    \فقره یک مورد 
%    \فقره یک مورد دیگر 
%    \فقره باز هم یک مورد دیگر
%\پایان{شمارش}

% برای نوشتن یک فرمول چند خطی که یک شماره بخوره. برای حالت‌های دیگر به سایت زیر مراجعه کنید:
% https://www.overleaf.com/learn/latex/Aligning_equations_with_amsmath
%\begin{align} 
%    \begin{split}
%        (x+y)^3 &= (x+y)^2(x+y) \\
%                &= (x^2+2xy+y^2)(x+y) \\
%                &= (x^3+2x^2y+xy^2) + (x^2y+2xy^2+y^3) \\
%                &= x^3+3x^2y+3xy^2+y^3
%    \end{split}                                 
%\end{align}

% نوشتن ماتریس و آرایه دوبعدی
%\begin{align}
%    A = 
%    \begin{bmatrix}
%        A_{11} & A_{21} \\
%        A_{21} & A_{22}
%    \end{bmatrix}
%\end{align}

% اضافه کردن تصویر با شرح. برای اضافه کردن تصویر بدون شرح، فقط دستور درج‌تصویر را استفاده کنید.
%\شروع{شکل}[h!]
%\تنظیم‌ازوسط
%  \درج‌تصویر[پهنا=0.5\پهنای‌متن]{img.jpg}
%\شرح{شرح تصویر}
%\برچسب{تصویر:۱}
%\پایان{شکل}
%\پایان{زیرمسائل}

%کد زیر را ببینید.
%%\lstinputlisting{source_filename.py}
%\begin{latin}
%\begin{lstlisting}
%# Import the modules
%import sys
%import random
%
%ans = True
%
%while ans:
%question = raw_input("Ask the magic 8 ball a question: (press enter to quit) ")
%
%answers = random.randint(1,8)
%
%if question == "":
%sys.exit()
%
%elif answers == 1:
%print "It is certain"
%
%elif answers == 2:
%print "Outlook good"
%
%elif answers == 3:
%print "You may rely on it"
%
%elif answers == 4:
%print "Ask again later"
%
%elif answers == 5:
%print "Concentrate and ask again"
%
%elif answers == 6:
%print "Reply hazy, try again"
%
%elif answers == 7:
%print "My reply is no"
%
%elif answers == 8:
%print "My sources say no"
%\end{lstlisting}
%\end{latin}
